\documentclass[11pt,a4paper]{article}
\usepackage[utf8]{inputenc}
\usepackage[T1]{fontenc}
\usepackage[italian]{babel}
\usepackage{geometry}
\geometry{top=2.5cm, bottom=2.5cm, left=2.5cm, right=2.5cm}
\usepackage{amsmath,amssymb,amsthm}
\usepackage{graphicx}
\usepackage{parskip}
\usepackage{hyperref}
\usepackage{enumitem}
\usepackage{microtype}
\usepackage{minted}

\title{\vspace{-1cm}Calcolo Numerico 2025-26\\[6pt]
\large \textbf{Homework 1 -- Zeri di funzione}}
\author{}
\date{}

\begin{document}
\maketitle
\section*{Consegna}
\noindent Zeri di funzione: utilizzare il metodo di bisezione, il metodo delle iterazioni di punto fisso e il metodo di Newton per il calcolo dello zero delle seguenti funzioni. Disegnare il grafico della funzione per localizzare lo zero e scegliere sia l'intervallo in cui applicare la bisezione che il punto iniziale per gli altri due algoritmi.

\bigskip

\begin{enumerate}[left=0pt,label=\textbf{\arabic*.}]
  \item \(f(x)=\ln(x+1)-x\). \quad ( \(g(x)=\ln(x+1)\) )
  \medskip

  \item \(f(x)=x^{2}-\cos(x)\). \quad ( \(g(x)=\sqrt{\cos(x)}\) )
  \medskip

  \item \(f(x)=\sin(x)-x^{2}\). \quad ( \(g(x)=2\sin(x)\) )
  \medskip

  \item \(f(x)=e^{x}-3x\). \quad ( \(g(x)=\tfrac{1}{3}e^{x}\) )
  \medskip
\end{enumerate}

\bigskip

\section*{Codici degli algoritmi}
\subsection*{Metodo di bisezione}

\begin{minted}[fontsize=\small, linenos, frame=single]{python}
    def metodo_bisezione(f, a, b, max_iter, epsilon):
    cond = True
    i = 0
    if f(a) * f(b) >= 0.0:
        return None, 0
    while cond and i < max_iter:
        c = (a + b) / 2
        if abs(f(c)) < epsilon:
            cond = False
        if cond and f(a) * f(c) < 0:
            b = c
        elif cond:
            a = c
        i = i + 1 
    return c, i
\end{minted}

\bigskip
\bigskip
\bigskip
\bigskip


\subsection*{Metodo delle iterazioni di punto fisso}

\begin{minted}[fontsize=\small, linenos, frame=single]{python}
    def metodo_punto_fisso(g, x0, max_iter, epsilon1, epsilon2):
    i = 0
    cond = True
    xk = x0
    while cond and i < max_iter:
        x = g(xk)
        if abs(x - xk) < epsilon1 or abs(f(x)) < epsilon2:
            cond = False
        xk = x
        i = i + 1
    return xk, i
\end{minted}


\subsection*{Metodo di Newton}

\begin{minted}[fontsize=\small, linenos, frame=single]{python}
   def metodo_newton(f, df, x0, max_iter, epsilon1, epsilon2):
    i = 0
    cond = True
    xk = x0
    if abs(df(xk)) < 1e-10:
        return None, 0
    while cond and i < max_iter:
        xk1 = xk - f(xk) / df(xk)
        if abs(xk1 - xk) < epsilon1 or abs(f(xk1)) < epsilon2:
            cond = False
        xk = xk1
        i = i + 1
    return xk, i
\end{minted}

\bigskip
\bigskip

\section*{Risultati delle esecuzioni}
\subsection*{Prima funzione:  \(f(x)=x^{2}-\cos(x)\). \quad ( \(g(x)=\sqrt{\cos(x)}\) )}


\begin{table}[h!]
\centering
\begin{tabular}{|c|c|c|c|c|c|}
\hline
\textbf{Metodo} & \textbf{Punto iniziale} & \textbf{Intervallo} & \textbf{Epsilon} & \textbf{Soluzione} & \textbf{Iterazioni}  \\ \hline
 Bisezione   & 0   &  [0,1]  & \(10^{-6}\) & 0.824131965637207 & 20   \\ \hline
 Punto Fisso & 0   &  [0,1]  & \(10^{-6}\) & 0.824132109264407 & 18   \\ \hline
 Newton      & 0   &  [0,1]  & \(10^{-6}\) &  &     \\ \hline
\end{tabular}
\end{table}



\begin{table}[h!]
\centering
\begin{tabular}{|c|c|c|c|c|c|}
\hline
\textbf{Metodo} & \textbf{Punto iniziale} & \textbf{Intervallo} & \textbf{Epsilon} & \textbf{Soluzione} & \textbf{Iterazioni}  \\ \hline
 Bisezione   & 0.5   &  [0,1]  & \(10^{-6}\) & 0.824131965637207 & 20   \\ \hline
 Punto Fisso & 0.5   &  [0,1]  & \(10^{-6}\) & 0.824132596440194 & 17   \\ \hline
 Newton      & 0.5   &  [0,1]  & \(10^{-6}\) & 0.824132312409912 & 4    \\ \hline
\end{tabular}
\end{table}



\begin{table}[h!]
\centering
\begin{tabular}{|c|c|c|c|c|c|}
\hline
\textbf{Metodo} & \textbf{Punto iniziale} & \textbf{Intervallo} & \textbf{Epsilon} & \textbf{Soluzione} & \textbf{Iterazioni}  \\ \hline
 Bisezione   & 0.8   &  [0,1]  & \(10^{-6}\) & 0.824131965637207 & 20   \\ \hline
 Punto Fisso & 0.8   &  [0,1]  & \(10^{-6}\) & 0.824132025062060 & 14   \\ \hline
 Newton      & 0.8   &  [0,1]  & \(10^{-6}\) & 0.824132376563292 & 2   \\ \hline
\end{tabular}
\end{table}





\begin{table}[h!]
\centering
\begin{tabular}{|c|c|c|c|c|c|}
\hline
\textbf{Metodo} & \textbf{Punto iniziale} & \textbf{Intervallo} & \textbf{Epsilon} & \textbf{Soluzione} & \textbf{Iterazioni}  \\ \hline
 Bisezione   & 0.8   &  [0,1]  & \(10^{-5}\) & 0.824134826660156 & 17   \\ \hline
 Punto Fisso & 0.8   &  [0,1]  & \(10^{-5}\) & 0.824125006492905 & 10   \\ \hline
 Newton      & 0.8   &  [0,1]  & \(10^{-5}\) & 0.824132376563292 & 2    \\ \hline
\end{tabular}
\end{table}



\begin{table}[h!]
\centering
\begin{tabular}{|c|c|c|c|c|c|}
\hline
\textbf{Metodo} & \textbf{Punto iniziale} & \textbf{Intervallo} & \textbf{Epsilon} & \textbf{Soluzione} & \textbf{Iterazioni}  \\ \hline
 Bisezione   & 0.8   &  [0,1]  & \(10^{-4}\) & 0.8240966796875   & 13   \\ \hline
 Punto Fisso & 0.8   &  [0,1]  & \(10^{-4}\) & 0.824215052092387 & 8   \\ \hline
 Newton      & 0.8   &  [0,1]  & \(10^{-4}\) & 0.824132376563292 & 2    \\ \hline
\end{tabular}
\end{table}




\begin{table}[h!]
\centering
\begin{tabular}{|c|c|c|c|c|c|}
\hline
\textbf{Metodo} & \textbf{Punto iniziale} & \textbf{Intervallo} & \textbf{Epsilon} & \textbf{Soluzione} & \textbf{Iterazioni}  \\ \hline
 Bisezione   & 0.8   &  [0,1]  & \(10^{-3}\) & 0.82421875 & 8   \\ \hline
 Punto Fisso & 0.8   &  [0,1]  & \(10^{-3}\) & 0.824549519096654 & 5   \\ \hline
 Newton      & 0.8   &  [0,1]  & \(10^{-3}\) & 0.824470434030340 & 1    \\ \hline
\end{tabular}
\end{table}




\subsection*{Metodo di Newton}

\begin{minted}[fontsize=\small, linenos, frame=single]{python}
   def metodo_newton(f, df, x0, max_iter, epsilon1, epsilon2):
    i = 0
    cond = True
    xk = x0
    if abs(df(xk)) < 1e-10:
        return None, 0
    while cond and i < max_iter:
        xk1 = xk - f(xk) / df(xk)
        if abs(xk1 - xk) < epsilon1 or abs(f(xk1)) < epsilon2:
            cond = False
        if cond:
            xk = xk1
            i = i + 1
    return xk, i
\end{minted}






\end{document}
